% Generated draft: requires mathematical and source review.
\documentclass[11pt]{article}
\usepackage[margin=25mm]{geometry}
\usepackage{fontspec}
\setmainfont{Latin Modern Roman}
\usepackage{amsmath,amssymb,mathtools,unicode-math}
\setmathfont{Latin Modern Math}
\usepackage{enumitem,needspace,xurl,hyperref}
\setlength{\parindent}{0pt}
\setlength{\parskip}{5pt}
\setlength{\emergencystretch}{3em}
\newcommand{\N}{\mathbb{N}}
\newcommand{\Z}{\mathbb{Z}}
\newcommand{\Q}{\mathbb{Q}}
\newcommand{\R}{\mathbb{R}}
\newcommand{\C}{\mathbb{C}}
\newcommand{\F}{\mathbb{F}}
\newcommand{\A}{\mathbb{A}}
\newcommand{\PP}{\mathbb{P}}
\newcommand{\E}{\mathbb{E}}
\newcommand{\eps}{\varepsilon}
\newcommand{\dd}{\,\mathrm{d}}
\begin{document}
\section*{UnsolvedMath ID: 700960}
\textbf{Source problem code:} catalogue entry
\subsection{Definitions and mathematical statement}
For a group \(H\), define its lower central series by
\[
\gamma_1(H)=H,\qquad
\gamma_{n+1}(H)=[\gamma_n(H),H]\quad(n\geq 1),
\]
where \([K,H]\) is the subgroup generated by the commutators \(k^{-1}h^{-1}kh\) with \(k\in K\) and \(h\in H\). The group \(H\) is nilpotent if \(\gamma_{c+1}(H)=\{1\}\) for some integer \(c\geq 1\). For a set \(X\), write \(F(X)\) for the free group on \(X\); a free nilpotent group has the form
\[
N=F(X)/\gamma_{c+1}(F(X)),\qquad c\geq 1.
\]
The source does not specify restrictions on the rank \(\lvert X\rvert\) or on \(c\).

Write \(T_A(N)\) provisionally for the \emph{tensor \(A\)-completion} referred to in A. G. Myasnikov and V. N. Remeslennikov, \emph{Siberian Math. J.} \textbf{35}, no. 5 (1994), 986--996. The supplied record does not define this construction, the nature of \(A\), or its admissibility hypotheses; this notation is a reference to that missing definition and does not specify a construction independently.

The requested assertion is that a free nilpotent group can have a non-nilpotent tensor \(A\)-completion. With admissibility understood to include all hypotheses of the cited construction, its provisional existential formulation is
\[
\exists\,(A,X,c)\ \text{admissible},\quad c\geq 1,\qquad
\forall n\geq 1,\quad
\gamma_{n+1}\!\left(T_A\!\left(F(X)/\gamma_{c+1}(F(X))\right)\right)\neq\{1\}.
\]
This formulation requires clarification of the admissible coefficient objects and input groups. In particular, the supplied wording does not determine whether \(A\) may be chosen or is intended to be fixed subject to further hypotheses.
\subsection{Short English statement}
The problem asks for a free nilpotent group whose tensor \(A\)-completion is not nilpotent. A precise unconditional formulation requires the missing definition of that completion and the hypotheses on \(A\). The supplied registry labels the problem open and records an explicitly unverified AI solution claim; neither assessment is independently verified here.
\subsection{Sources}
\begin{itemize}
\item Supplied local UnsolvedMath registry record, permanent ID \(700960\), exact source problem code \texttt{catalogue entry}, title \emph{Kourovka Notebook Problem 13.42}; upstream revision \texttt{d371ee917ac9647fbb06cd6fc529deec6f3dda82}.
\item \emph{The Kourovka Notebook (Unsolved Problems in Group Theory)}, original problem 13.42, 13th issue (1995); attributed to A. G. Myasnikov and V. N. Remeslennikov. The record cites the 21st issue edition (2026), \texttt{arXiv:1401.0300v48}. The supplied collection description names E. I. Khukhro and V. D. Mazurov as editors.
\item A. G. Myasnikov and V. N. Remeslennikov, \emph{Siberian Math. J.} \textbf{35}, no. 5 (1994), 986--996, cited by the source for tensor \(A\)-completion. No article title or definition is supplied.
\item The supplied discussion records: [Unverified AI claim] an example is produced in M. Kinyon and J. Urban, \emph{Preprint} (2026), \url{https://kourovkanotebookorg.wordpress.com/wp-content/uploads/2026/09/kourovka-mathematics.pdf}. This is explicitly an unverified AI solution claim, not an established solution in this draft.
\item The supplied status review is dated \texttt{2026-10-16}, gives canonical status \texttt{open} and evidence \texttt{unrefereed\_claim}, and reports a \(\#\) marking for an unverified AI solution claim in the October 2026 21st issue. This assessment is supplied dated curation, not independently verified.
\end{itemize}
\end{document}
